\section{Validator set and decentralisation roadmap}
\label{sec:validators}

\subsection{The launch set, stated plainly}

At mainnet genesis Konstellation will be validated by \textbf{five to ten nodes
operated by the founding team and its foundation}, spread across multiple cloud
and bare-metal providers, autonomous systems and jurisdictions. That is a
\textbf{permissioned validator set}, and this document says so rather than
describing it as decentralised (\dref{D7}).

The reasons are the ones behind every other conservative choice here. The chain
holds real user funds from block one. A set of five to ten operators who share a
paging rota, a runbook and an incident channel can execute a security-patch
binary swap in under an hour (\S\ref{sec:neverfork}); a set of forty
independent operators, on the evidence of the 2026 incidents, cannot do it in
twenty. The launch set is small because the response window has to be.

What the launch set is \emph{not}: it is not a proof-of-authority chain. The
staking module is the standard Cosmos SDK delegated-proof-of-stake module with
the parameters in \S\ref{sec:staking}; validators bond \kash{}, delegators may
delegate to them, rewards and slashing are on-chain and pro-rata, and the
active set is selected by bonded stake up to the cap of 30. The 120\,M \kash{}
validator-bootstrap allocation (\S\ref{sec:allocation}) is what is bonded at
genesis, distributed across the launch validators via their genesis
transactions.

\todo{State the mechanism by which the set is permissioned at launch. \dref{D7}
records the posture, not the enforcement: whether the \code{MsgCreateValidator}
message type is disabled through the circuit breaker (\dref{D14}) until the set
opens, or whether ``permissioned'' simply means that the launch validators
hold the bonded genesis stake and outside validators are not actively admitted.
The two read very differently to an operator.}

\subsection{What holds the launch set accountable}

Because the set is small and foundation-operated, the operational controls
matter more, not less, and they are the same controls an external operator
would be held to:

\begin{itemize}
\item \textbf{Slashing applies to the foundation's validators} exactly as to
  anyone else: 5\,\% and a permanent tombstone for double-signing, 0.01\,\% and
  jail for downtime (\dref{D10}).
\item \textbf{Key isolation.} Each validator's consensus key is held by a
  threshold signer (Horcrux, three or more regions) or an HSM-backed signer, never
  on a machine that also serves public traffic. Running two nodes with the same
  consensus key is forbidden under all circumstances, including migrations and
  failover (ENGINEERING \S2.7, \S9.2).
\item \textbf{Reproducible binaries.} No validator runs a binary it built itself;
  every release is built in CI from a signed tag, twice, and published with
  checksums that operators verify before staging (ENGINEERING \S2.6).
\item \textbf{Provider and jurisdiction diversity} across the set, so that a
  single provider outage or a single jurisdiction's action cannot halt the chain
  on its own (ENGINEERING \S9.2).
\item \textbf{Every emergency power is on-chain and logged.} The circuit breaker,
  the compliance lists and the IBC rate limits all act through on-chain messages
  that emit events; governance can override each of them
  (\S\ref{sec:rails}, \S\ref{sec:compliance}).
\end{itemize}

\subsection{Roadmap for opening the set}
\label{sec:roadmap}

The direction is decided: validators are added over time as the network
decentralises, and this roadmap is the published form of that commitment
(\dref{D7}). The \emph{stages} below follow from the launch sequence
(\S\ref{sec:launch}) and decisions already on record. The \emph{dates and
numeric triggers} are not yet decided and are marked as such; they will be
filled in a later release of this document, not guessed here.

\begin{description}[style=nextline]
\item[Stage 0 --- in-house testnet (launch-sequence phase 5).]
  Five in-house validators on \code{testnet-1} for a minimum of six to eight
  weeks, including a state-breaking upgrade drill, a chaos test (40\,\% of
  validators killed mid-block) and a halt-and-restart drill.

\item[Stage 1 --- external operators on testnet (phase 6).]
  Three to five external operators are invited onto \code{testnet-1}. Their job
  is to find the gaps in the validator documentation that in-house engineers
  cannot see; the documentation is fixed before mainnet.

\item[Stage 2 --- mainnet genesis (phases 8--9).]
  Five to ten foundation-operated validators, permissioned, under a value
  ceiling (\S\ref{sec:bridges}). Soak for a month.

\item[Stage 3 --- first external mainnet validators.]
  Operators who ran \code{testnet-1} in Stage 1 are the natural first
  candidates. \todo{Decide the trigger for Stage 3 (a clean soak? the audit
  report published? a calendar date?), the number of operators admitted, and
  the admission criteria (testnet participation, key-management requirements,
  jurisdiction). None is recorded.}

\item[Stage 4 --- permissionless entry within the cap.]
  Any operator meeting the on-chain requirements (minimum self-delegation,
  5\,\% minimum commission) may create a validator and compete for one of the
  30 active slots by bonded stake. \todo{Decide when Stage 4 begins and whether
  a chain-wide minimum self-delegation floor is wanted (it requires a custom
  ante decorator; ENGINEERING \S8).}

\item[Stage 5 --- raising the cap and lengthening governance.]
  \code{max\_validators} (30) and the 3-day voting period are both governance
  parameters. The intent recorded with \dref{D11} is to lengthen the voting period
  as the set decentralises; the intent recorded with \dref{D10} is that 30 is a
  launch cap, not a permanent one. \todo{Decide target set sizes and voting
  periods for Stage 5, and whether they are tied to a share of stake held
  outside the foundation.}
\end{description}

Two properties of the economics are designed to make opening the set work
without further intervention. First, the issuance curve
(\S\ref{sec:issuance}) pays $\approx$11.5\,\% at the genesis bonding level and
falls as stake arrives, so outside delegation is pulled in when the set is
smallest and the pull fades as it grows. Second, the yield is chain-wide and
pro-rata, so adding validators changes who earns it, not how much is earned;
there is no economic reason for the foundation to keep the set small.

\subsection{What ``decentralised enough'' will mean}

\todo{Define the metrics this roadmap will be measured against in later
releases (for example: share of bonded stake held outside the foundation,
Nakamoto coefficient of the active set, number of distinct operating entities
and jurisdictions). No metric has been chosen.}
